\einheitskopf{Blatt 0 · Das kannst du schon}\label{b0}

\begin{aufgabe}{Brüche kürzen und auf Hundertstel erweitern. Kürze so weit wie möglich.}
\begin{teile}
\teil $\frac{4}{10} = \leerfeld$
\teil $\frac{6}{9} = \leerfeld$
\teil $\frac{15}{40} = \leerfeld$
\end{teile}
Erweitere auf den Nenner 100.
\begin{teile}
\teil $\frac{7}{20} = \frac{\;\;}{100}$ \quad Zähler: \leerfeld
\end{teile}
Kreuze an: Welcher Bruch lässt sich \emph{nicht} auf den Nenner 100 erweitern?
\begin{teile}
\teil \kreuz{$\frac{1}{4}$}\kreuz{$\frac{7}{8}$}\kreuz{$\frac{3}{10}$}
\end{teile}
\end{aufgabe}

\begin{aufgabe}{Brüche und Dezimalzahlen. Schreibe als Dezimalzahl.}
\begin{teilezwei}
\tz $\frac{1}{2} = \leerfeld$ & \tz $\frac{1}{4} = \leerfeld$ \\
\end{teilezwei}
Schreibe als Bruch und kürze so weit wie möglich.
\begin{teile}
\teil $0{,}4 = \leerfeld$
\end{teile}
Schreibe als Dezimalzahl.
\begin{teilezwei}
\tz $\frac{3}{5} = \leerfeld$ & \tz $\frac{1}{8} = \leerfeld$ \\
\end{teilezwei}
Schreibe als Bruch mit dem Nenner 100.
\begin{teile}
\teil $0{,}07 = \frac{\;\;}{100}$ \quad Zähler: \leerfeld
\end{teile}
\end{aufgabe}

\begin{aufgabe}{Auf eine Stelle nach dem Komma runden. Runde auf eine Stelle nach dem Komma.}
\begin{teilezwei}
\tz $3{,}24 \approx \leerfeld$ & \tz $7{,}68 \approx \leerfeld$ \\
\tz $12{,}349 \approx \leerfeld$ & \tz $9{,}96 \approx \leerfeld$ \\
\tz $5{,}55 \approx \leerfeld$ & \tz $28{,}04 \approx \leerfeld$ \\
\end{teilezwei}
\end{aufgabe}

\begin{aufgabe}{Durch 100 teilen und mit Dezimalzahlen malnehmen. Teile durch 100.}
\begin{teile}
\teil $400 : 100 = \leerfeld$
\teil $35 : 100 = \leerfeld$
\teil $7{,}50 : 100 = \leerfeld$
\end{teile}
Berechne.
\begin{teile}
\teil $0{,}4 \cdot 60 = \leerfeld$
\teil $0{,}04 \cdot 60 = \leerfeld$
\teil $1{,}2 \cdot 50 = \leerfeld$
\end{teile}
\end{aufgabe}

\begin{aufgabe}{In der Tabelle hoch- und runterrechnen. Rechne im Schema: erst auf 1 herunter, dann wieder hoch. Schreibe die Rechenschritte an die Pfeile.}
Beispiel:
\begin{dreisatz}{Hefte}{Preis}
\dsz{4}{$6{,}00\,€$}
\dsp{:4}{:4}
\dsz{1}{$1{,}50\,€$}
\dsp{\cdot 6}{\cdot 6}
\dsz{6}{$9{,}00\,€$}
\end{dreisatz}
\begin{teile}
\teil 5 Schlüsselanhänger kosten $4{,}00\,€$. Was kosten 8 Anhänger?
\begin{dreisatz}{Anhänger}{Preis}
\dsz{5}{$4{,}00\,€$}
\dsp{\phantom{:00}}{\phantom{:00}}
\dsz{1}{\dsleer}
\dsp{\phantom{:00}}{\phantom{:00}}
\dsz{8}{\dsleer}
\end{dreisatz}
\teil 6 Brötchen kosten $2{,}10\,€$. Was kosten 4 Brötchen?
\begin{dreisatz}{Brötchen}{Preis}
\dsz{6}{$2{,}10\,€$}
\dsp{\phantom{:00}}{\phantom{:00}}
\dsz{1}{\dsleer}
\dsp{\phantom{:00}}{\phantom{:00}}
\dsz{4}{\dsleer}
\end{dreisatz}
\teil 3\,m Geschenkband kosten $7{,}50\,€$. Was kosten 7\,m?
\begin{dreisatz}{Band in m}{Preis}
\dsz{3}{$7{,}50\,€$}
\dsp{\phantom{:00}}{\phantom{:00}}
\dsz{1}{\dsleer}
\dsp{\phantom{:00}}{\phantom{:00}}
\dsz{7}{\dsleer}
\end{dreisatz}
\teil 4 Lose kosten $3{,}00\,€$. Wie viele Lose bekommt man für $12{,}00\,€$?
\begin{dreisatz}{Lose}{Preis}
\dsz{4}{$3{,}00\,€$}
\dsp{\phantom{:00}}{\phantom{:00}}
\dsz{1}{\dsleer}
\dsp{\phantom{:00}}{\phantom{:00}}
\dsz{\dsleer}{$12{,}00\,€$}
\end{dreisatz}
\end{teile}
\end{aufgabe}

\begin{aufgabe}{Bruchteil einer Größe berechnen. Berechne und schreibe die Einheit dazu.}
\begin{teile}
\teil $\frac{1}{2}$ von $80\,€ = \leerfeld$
\teil $\frac{1}{4}$ von $60\,$kg $= \leerfeld$
\teil $\frac{3}{4}$ von $24\,$m $= \leerfeld$
\teil $\frac{2}{5}$ von $45\,€ = \leerfeld$
\teil $\frac{1}{10}$ von $7\,€ = \leerfeld$
\end{teile}
\end{aufgabe}

\begin{aufgabe}{Fehler finden. Mia soll $\frac{3}{8}$ als Dezimalzahl schreiben. Sie schreibt: $\frac{3}{8} = 0{,}38$. Finde den Fehler, benenne ihn und korrigiere die Rechnung.}
\begin{teile}
\teil Fehler: \dsleer\dsleer\dsleer\dsleer\ \ richtig: \leerfeld
\teil Schreibe $\frac{9}{20}$ als Dezimalzahl. \leerfeld
\end{teile}
\end{aufgabe}
